\chapter{Engineering at civilization scale}
\label{vol12:ch01}

\epigraph{The bridge is engineering; the network of bridges, the policies that built them, the populations they connect and displace, the maintenance that keeps them up and the budget that does not, is also engineering.}{}

\chapmeta{Half-life: medium-life. Archetypes: scaling, ethics. Exercise target: 15. Project track: hybrid.}

\noindent
This volume turns the engineering discipline from the artifact to the civilizational scale at which engineered systems constitute the environment human populations live in. The chapter's claim is that the larger scale is still engineering: the same physical reality; the same mathematical discipline; the same failure literacy; the same design responsibility. What changes is the actors, the time horizons, the integration across systems, and the political-economic context within which the engineering work is done. The chapter develops what these changes require of the engineer.

\section{What scale changes}\pathtag{core}

A bridge is engineering. The network of bridges in a metropolitan transportation system is engineering at a different scale. The two are continuous in the underlying disciplines (mechanics; materials; construction; maintenance) and discontinuous in several respects that this volume develops:

\textbf{Number of actors.} A bridge is designed by a small engineering team for a single client. The network is the cumulative product of decades of design decisions by hundreds of engineering teams for many clients (municipalities; states; federal authorities; private developers; community boards). The integration of the decisions is itself a design problem; the integration is rarely the responsibility of any single party.

\textbf{Time horizons.} A bridge is designed for a service life of 50-100 years. The network exists across centuries; the policy frameworks that govern it span decades; the political cycles that fund it are 2-6 years. The mismatch among these horizons is structural; the engineering work at network scale includes managing the mismatch.

\textbf{Integration across systems.} A bridge serves transportation; the network serves transportation, land use, economic geography, environmental footprint, climate response, equity. The bridge's design considers structural and traffic factors; the network's design considers all of these and more, with the trade-offs between them as the substantive engineering work.

\textbf{Distributive consequences.} A bridge serves the users who cross it. The network serves some populations and harms others: the populations it connects benefit; the populations it does not connect are excluded; the populations whose neighbourhoods are bisected or demolished bear the cost. The distributive consequences are not incidental to the engineering; they are the engineering's substantive output.

\textbf{Maintenance over generations.} A bridge requires maintenance over its lifecycle (Volume X Chapter 14). The network requires maintenance across generations of operators, funders, and users; the maintenance discipline must persist across the personnel changes that the time horizon implies. The cumulative maintenance is the network's principal lifecycle cost.

\textbf{Failure modes at scale.} A bridge failure is local. A network failure (a hurricane that takes out the bridge plus its alternatives; a cyberattack on the traffic-management system; a fuel-supply disruption that disables the trucking network) is regional or national. The failure modes scale; the engineering responses must scale with them.

\textbf{Political-economic context.} A bridge is built with the resources the client commands. The network exists within a political-economic framework that allocates resources across many systems and many decades; the engineering work at network scale includes engaging with the framework rather than treating it as external to engineering.

The scale change is not merely quantitative; it is qualitative. The engineering at civilization scale is engineering in the substantive sense of Volume I's thesis ("Engineering is disciplined intervention"), but the discipline includes commitments and responsibilities that artifact-scale engineering can sometimes avoid. The volume's chapters develop the specific commitments by domain.

\begin{principle}[The scale is the discipline]
Engineering at civilization scale is not a different discipline from artifact-scale engineering; it is the same discipline applied to a scale where the actors, time horizons, integrations, distributions, and political-economic contexts are part of the engineering's substantive content. The engineer who treats the larger scale as outside the engineering profession's responsibility has misread what engineering at this scale requires.
\end{principle}

\section{The unit of analysis: from artifact to infrastructure to system-of-systems}\pathtag{core}

The engineering literature's units of analysis nest:

\textbf{Artifact.} A discrete engineered object: a bridge, a building, a vehicle, a circuit, a software application, a medical device. The artifact has well-defined boundaries; its design and analysis are the subject of most of this book's earlier volumes. The artifact serves a specific function; its requirements, performance, and failure modes are characterised at the artifact level.

\textbf{Infrastructure.} A network of artifacts that collectively serve a function at population scale. The road network; the electrical grid; the water and wastewater systems; the communications network; the air-traffic system; the hospital system. Infrastructure has many constituent artifacts; the artifacts interact; the interactions produce behaviours the individual artifacts do not. Infrastructure has institutional owners, operators, regulators, and users; the institutional structure is part of the infrastructure's existence.

\textbf{System-of-systems.} An aggregation of infrastructures that interact: the energy-transportation-information system that constitutes modern logistics; the energy-water-food system that constitutes the food supply; the healthcare-water-housing system that constitutes public health. The system-of-systems has emergent properties from the interactions of its constituent infrastructures; the failures cascade across infrastructures (Volume X Chapter 7 cascading failure); the design of any single infrastructure has consequences for the other infrastructures it interacts with.

The engineering disciplines that operate at each level differ:

\textbf{Artifact-level engineering.} The disciplines this book has developed in Volumes I-XI: mechanics, materials, energy, control, software, design. The output is an artifact that meets its requirements within its operating envelope.

\textbf{Infrastructure-level engineering.} The disciplines of network design, capacity planning, operations management, regulatory engagement, maintenance organisation, demand forecasting, asset management. The output is an infrastructure that serves a population over decades. The discipline draws on artifact-level engineering for the components and adds the infrastructure-level concerns of network behaviour, institutional structure, and lifecycle management.

\textbf{System-of-systems engineering.} The disciplines of inter-infrastructure coordination, cross-system failure analysis, regulatory coherence across sectors, long-term planning across institutional boundaries, integrated investment decisions. The output is a system-of-systems that functions coherently despite being assembled from independent infrastructures with independent governance. The discipline is the least mature of the three; the institutional structures for it are still evolving.

The civilization-scale engineering of this volume operates predominantly at the infrastructure and system-of-systems levels. The artifact-level engineering of Volumes I-XI is the prerequisite; the civilization-scale work is what those prerequisites enable.

The discipline of working at the larger scales includes:

\textbf{Decomposition that preserves emergent behaviour.} Civilization-scale systems can be decomposed for analysis (the road network into its bridges and roadways; the grid into its generation, transmission, and distribution; the food system into its production, processing, and distribution). The decomposition is useful but cannot lose the emergent behaviour that arises from the integration. The discipline includes both the decomposition and the maintenance of the integrated perspective.

\textbf{Multi-stakeholder engineering.} The engineering work involves stakeholders whose interests differ: users versus operators versus regulators versus affected communities. The work includes engaging with each, surfacing the conflicts, designing the engineering response that addresses the conflicts substantively rather than ignoring them.

\textbf{Long-horizon design.} The design horizons are decades to centuries. The design must accommodate uncertainty about future conditions: climate change; demographic shifts; technological change; political evolution. The discipline includes designing for robustness across plausible futures rather than for optimisation against a single forecast.

\textbf{Inter-infrastructure dependency.} The engineering work includes characterising the dependencies between infrastructures: the grid depends on the gas supply; the gas supply depends on the trucking network; the trucking network depends on the road network; the road network depends on the construction industry; the construction industry depends on the materials supply; the materials supply depends on the grid. The dependencies are circular; the design must address the circular structure.

\begin{archetype}[scaling]
Civilization-scale engineering is the scaling archetype's planetary expression. The engineering work at the artifact level operates within bounded scope; the engineering work at the civilization level operates within a system whose extent is the planet's population and whose timescale is the human civilisation's continuation. The scaling produces qualitatively different engineering, not just more of the same.
\end{archetype}

\section{Time horizons engineers have to learn}\pathtag{core}

Civilization-scale engineering operates across time horizons that artifact-scale engineering does not normally engage:

\textbf{Decadal.} The horizon of major infrastructure investments: a new airport; a transmission line; a wastewater treatment plant. The decisions made today commit the population to the resulting infrastructure for decades.

\textbf{Multi-generational.} The horizon of foundational infrastructure: the major water systems built in the late 19th century; the road networks of the mid-20th century; the public buildings of every century. The infrastructure outlives its designers, its builders, and often the institutions that commissioned it.

\textbf{Multi-century.} The horizon of architectural and civil works whose use is intended across centuries: cathedrals; aqueducts; some bridges; some buildings. The design discipline includes provisions for maintenance across the centuries; the institutional commitment is the design's substantive prerequisite.

\textbf{Climate-relevant.} The horizon of climate change consequences: the carbon emitted now affects atmospheric concentrations for centuries; the sea-level rise responses to current emissions extend over millennia. The engineering decisions made now have consequences across time horizons longer than any human life.

\textbf{Radioactive-waste-relevant.} The horizon of nuclear waste management: the waste from current operations remains hazardous for thousands of years; the engineering institutions for its long-term storage must function across timescales longer than any institution has yet functioned. The engineering discipline includes designing the storage and the institutions for the storage; both are part of the engineering work.

The engineer at civilization scale must work in time horizons their personal experience and intuition do not support. The discipline includes:

\textbf{Cohort thinking.} The decisions are made by one cohort of engineers and operate across many subsequent cohorts. The cohort thinking discipline includes documenting the decisions in forms that subsequent cohorts can understand and act on, building in flexibility for conditions the current cohort cannot foresee, and accepting the responsibility for decisions whose consequences the current cohort will not see.

\textbf{Discounting reluctance.} The financial discipline of discounting future costs at current interest rates produces present-value calculations that systematically under-weight long-term consequences. The civilization-scale engineering discipline includes resistance to discounting when the consequences are not financial but environmental or human-life: the future generations have no representation in current market interest rates; their interests must be represented through engineering and policy judgment.

\textbf{Institutional design for persistence.} The engineering work includes the design of the institutions that will maintain the engineering across time. The institutions are themselves engineered systems with their own failure modes (Volume X Chapter 10 on organisational failure); the design includes anti-fragility against the institutional failure modes.

\textbf{Documentation for posterity.} The engineering decisions are documented in forms accessible across the time horizons. The Yucca Mountain nuclear waste repository programme commissioned studies of how to communicate the hazard to populations 10{,}000 years from now; the studies produced symbol systems and architectural forms intended to convey "do not dig here" across cultural and linguistic discontinuities. The discipline is in early stages; its substantive value is for engineering decisions whose consequences extend beyond the documentation systems currently in place.

\textbf{Reversibility preference.} Where commitments can be reversible without substantial cost, the discipline prefers reversible commitments. The reversibility preserves the option to revise decisions when subsequent information becomes available; irreversibility commits the future to the current decision without the future's ability to revise it.

\section{The engineer as one of many actors}\pathtag{standard}

Civilization-scale engineering is not the work of any single engineer or even of any single engineering team. The work is the product of many actors whose roles include and extend beyond engineering:

\textbf{Engineers.} The technical work the previous volumes have developed: design, analysis, prototyping, testing, sign-off. The engineer's role is necessary; it is not sufficient.

\textbf{Project managers.} The coordination of multi-team engineering work: schedule, budget, scope, stakeholder management. The project-management discipline is itself engineering in the wide sense; its quality affects whether the technical work delivers.

\textbf{Architects and urban planners.} The design of spatial arrangements at building and city scale. The discipline overlaps with engineering and has its own framing, training, and professional structures.

\textbf{Clients and procurers.} The parties who fund the engineering work and specify (often inadequately) what they want. The client's substance shapes the engineering's scope; the engineer's engagement with the client is itself part of the engineering work.

\textbf{Regulators.} The bodies that constrain the engineering's scope through standards, permits, certifications, and ongoing oversight (Volume X Chapter 11; Volume XI Chapter 9). The regulator's engagement with the engineering work is structured; the engineering work includes the engagement.

\textbf{Operators and maintainers.} The parties who run the engineered systems after they are built: utility operators, traffic managers, hospital staff, telecommunications operators, port operators. The operational discipline is the engineered system's continued reality; the engineering work includes designing for the operators' actual capabilities and incentives.

\textbf{Users.} The populations who interact with the engineered systems: the drivers, the patients, the consumers, the citizens. The user research (Volume XI Chapter 3) develops the discipline; civilization-scale engineering extends it to the broader populations the systems serve and exclude.

\textbf{Affected non-users.} The populations who do not use the engineered system but who are affected by it: the residents whose neighbourhoods are bisected by the highway; the downstream populations whose water quality is affected by the upstream plant; the future generations exposed to current decisions' climate consequences. The engineering work at civilization scale includes engaging with these populations whose interests artifact-scale engineering can often ignore.

\textbf{Politicians.} The elected and appointed officials who allocate resources and set policy. The engineering work engages with politicians as the substantive decision-makers for civilization-scale projects; the engagement requires translating engineering judgment into terms politicians can act on without losing the engineering substance.

\textbf{Funders and investors.} The parties who provide the capital for the engineering work: governments, private investors, multilateral institutions, philanthropies. The funding's structure shapes what gets built and how; the engineering work includes engaging with the funders.

\textbf{Other engineers.} The engineers in adjacent fields, in regulatory bodies, in academic positions, in international standards bodies. The engineering profession is the collective whose judgment shapes the standards within which the engineering work is done.

The engineer at civilization scale is one of these many actors. The engineering work is substantive (the technical decisions matter; the engineering's quality determines whether the system works); the engineering work is also embedded in the actor network that produces the system. The discipline of working within the network includes:

\textbf{Knowing one's role.} The engineer is not the project manager, not the architect, not the regulator, not the politician. The engineer's contribution is technical; the contribution is enriched by understanding the other actors' roles but is not improved by attempting to play them all.

\textbf{Engaging substantively with adjacent disciplines.} The engineering's quality depends on the engineer's engagement with the adjacent disciplines: the project management's discipline; the architectural and urban-planning discipline; the regulatory discipline; the operational discipline. The engagement is the substance of working at civilization scale.

\textbf{Communicating engineering to non-engineers.} The substantive engineering work must be communicated to actors who are not engineers: the politicians who allocate; the funders who finance; the affected populations who bear the consequences; the public whose collective decisions shape the framework. The communication is the engineering profession's interface with the broader society; its quality shapes whether the engineering work serves the society or is captured by interests less aligned with the public.

\textbf{Resisting both excessive deference and excessive autonomy.} The engineer who defers excessively to non-engineering actors becomes a tool whose judgment is not exercised; the engineer who acts excessively autonomously bypasses the legitimate decision-makers who have authority the engineer lacks. The discipline is the balance.

\section{The thesis revisited at civilization scale}\pathtag{core}

The thesis of this book, pinned at the start of Volume I and recurring throughout:

\begin{quote}
\emph{Engineering is the discipline by which measured reality becomes reliable intervention under constraint, failure, scale, and responsibility. Engineering is not applied science. Engineering is disciplined intervention.}
\end{quote}

At civilization scale, the thesis's elements take on extended meaning:

\textbf{Measured reality.} At artifact scale, the measurement is of forces, temperatures, voltages. At civilization scale, the measurement includes the population's needs, the climate's trajectory, the economy's structure, the ecological state. The measurement disciplines that produce the larger-scale data (census; epidemiology; remote sensing; economic indicators; environmental monitoring) are the engineering profession's foundation at this scale.

\textbf{Reliable intervention.} At artifact scale, the intervention is the artifact's function over its service life. At civilization scale, the intervention is the infrastructure or the system-of-systems serving the population over its multi-generational existence. The reliability discipline scales accordingly: the reliability of an electrical grid is the population's expectation of always having power; the reliability of a transportation network is the population's expectation of always being able to move.

\textbf{Constraint.} The constraints at civilization scale include the constraints familiar at artifact scale (physics; materials; economics) and additional constraints specific to the larger scale: the planet's biophysical boundaries; the population's tolerance for disruption; the political framework's continued legitimacy; the cumulative environmental cost. The engineer at this scale operates within constraints richer than at artifact scale.

\textbf{Failure.} The failure modes at civilization scale include the failure modes familiar at artifact scale (Volume X) and additional modes specific to the larger scale: civilizational collapse; ecological collapse; institutional collapse; political collapse. The failure literacy of Volume X is the foundation; the civilization-scale failure literacy extends it.

\textbf{Scale.} The scale itself is part of the thesis. The engineering at small scale serves bounded populations; the engineering at civilization scale serves humanity's continued capacity to inhabit the planet. The scaling produces qualitatively different engineering; the discipline includes recognising the qualitative differences and engaging with them substantively.

\textbf{Responsibility.} The responsibility at artifact scale is to the artifact's users and to the engineering profession. The responsibility at civilization scale extends to the population the system serves, to the affected non-users, to the future generations, to the planet's other inhabitants. The engineer's professional commitment includes the responsibility's broader scope.

\textbf{Disciplined intervention.} The discipline is what the previous volumes developed; the intervention is the engineering work. At civilization scale, both terms have extended meaning: the discipline includes engagement with the political economy and the institutional framework; the intervention includes the cumulative effect on the planet's human and ecological future.

\begin{principle}[The thesis at civilization scale]
The engineering profession's discipline of measured-reality-to-reliable-intervention operates at scales from the single artifact to the civilization that the artifacts collectively constitute. The discipline is continuous; the responsibilities scale; the engineer who exercises the discipline at artifact scale and ignores it at civilization scale has accepted the smaller responsibility without the larger one. The volume's chapters develop the larger responsibility's substance.
\end{principle}

\section{Failure: civilization-scale projects that did not understand they were civilization-scale}\pathtag{core}

The failure mode this section addresses is the project that was approached as if it were artifact scale or infrastructure scale when its actual scale and consequences were civilizational. The mismatch produces designs that work at the assumed scale and fail at the actual scale.

The recurring patterns:

\textbf{The Aral Sea collapse.} The Soviet irrigation diversion of the Amu Darya and Syr Darya rivers in the 1960s was approached as an agricultural-engineering project: divert water for cotton production; achieve agricultural output targets. The project's civilizational consequences included the desiccation of the Aral Sea (a substantial inland sea reduced to approximately 10\% of its original volume by 2010); the collapse of the regional fishing economy; the public-health consequences from dust storms carrying pesticide residues from the exposed seabed; the geopolitical tensions across the four post-Soviet states that share the rivers. The engineering's artifact-scale success (the diversion worked; cotton was produced) was the civilizational failure (the regional ecosystem and economy collapsed; the populations dependent on the sea were displaced; the political-economic consequences continue).

\textbf{The American interstate highway system.} The Federal-Aid Highway Act of 1956 commissioned a national interstate highway network as an infrastructure-engineering project. The project's civilizational consequences included the restructuring of American urban form (suburbanisation; central-city decline; automobile dependence); the racial and economic geography of postwar America (the highways were repeatedly routed through Black and immigrant neighbourhoods, displacing populations whose political representation was inadequate to resist); the long-term lock-in to automobile-based transportation that constrains the country's response to climate change; the cumulative environmental footprint. The engineering's infrastructure-scale success (the network was built; the design served its functional requirements) was a civilization-scale outcome with consequences the original engineering work did not address.

\textbf{The Bhopal disaster.} (Volume X Chapter 10.) Union Carbide's pesticide plant in Bhopal was approached as a chemical-process-engineering project. The plant's location in a populated area, the parent corporation's progressive disinvestment, the regulatory framework that did not enforce safety requirements, and the eventual catastrophic release on 2-3 December 1984 reflected the project's actual civilization-scale context. The engineering's process-engineering scope was inadequate to the civilization-scale consequences.

\textbf{The Three Gorges Dam.} The Three Gorges Dam in China was approached as a hydroelectric-engineering project with substantial flood-control benefits. The dam's civilizational consequences include the displacement of approximately 1.3 million people; the loss of substantial cultural heritage upstream; the ecological consequences for the Yangtze River system; the sediment-management consequences downstream including coastal-erosion effects in the Yangtze delta. The engineering's hydroelectric scope produced the dam; the civilizational consequences include displacement, ecological transformation, and downstream geomorphological changes.

\textbf{The Internet.} The early packet-switched network designs of the 1970s and 1980s were approached as communications-engineering projects to enable robust inter-computer communication. The civilizational consequences include the restructuring of commerce, media, social organisation, political process, and human cognition over the subsequent decades. The original engineering's scope was technical (protocols; routing; addressing); the civilizational consequences include effects the engineering community has been substantively engaged in addressing for decades and has not fully resolved.

\textbf{Antibiotic discovery and distribution.} The discovery of penicillin and the subsequent antibiotic revolution were approached as medical-engineering and pharmaceutical-engineering projects. The civilizational consequence include the extension of human life expectancy by decades and the emergence of antimicrobial resistance as a planetary-scale problem (a "civilization-scale failure mode" that the original engineering did not anticipate). The engineering response to antimicrobial resistance is itself civilization-scale work in progress.

The recurring lessons from these cases:

\textbf{The project's actual scale exceeds its assumed scale.} The engineering teams approached projects at the scale of their immediate technical scope; the projects' actual consequences operated at scales the technical scope did not address. The engineering response includes the discipline of asking, at project inception, what the project's actual scale is and whether the engineering scope is adequate to it.

\textbf{The civilization-scale consequences fall on populations without engineering representation.} The displaced populations, the future generations, the affected ecosystems, the downstream beneficiaries and victims, the populations bearing the unintended consequences: these typically lack institutional representation at the project-decision stage. The engineering profession's discipline at civilization scale includes the substantive engagement with the under-represented populations' interests.

\textbf{The institutional structures lag the engineering capability.} The engineering capability to produce civilization-scale interventions consistently exceeds the institutional structures' capacity to govern them. The institutional structures evolve in response (Volume X Chapter 11 on regulatory failure documents some of the evolution); the evolution lags the engineering capability by decades; the consequences accumulate in the gap.

\textbf{The engineering profession's responsibility includes the institutional gap.} The engineering profession can advocate for institutional structures adequate to the engineering capability; can refuse to participate in projects whose institutional governance is inadequate; can contribute its technical expertise to the institutional design that the profession is implicitly the principal beneficiary and the principal contributor to. The professional commitment is part of the engineering at civilization scale.

\begin{principle}[Civilization-scale engineering requires civilization-scale engagement]
The engineering profession's discipline includes the discipline of recognising when projects are civilization-scale in their consequences and engaging the engineering work to the scale of the consequences. Engineering scope that addresses only the artifact-scale technical work when the project's scale is civilizational produces results that satisfy the artifact-scale requirements and produce civilizational failures the requirements did not address.
\end{principle}

\begin{archetype}[ethics]
Civilization-scale engineering is the ethics archetype's planetary expression. The engineering decisions affect populations whose interests the engineering team must engage substantively; the engineering responsibilities extend across actors, time horizons, and consequences the artifact-scale work does not. The discipline of the engineer at this scale is the discipline of professional responsibility scaled to the scope of the engineering work's actual consequences.
\end{archetype}

\begin{project}[Trace a civilization-scale project]
Select one civilization-scale engineering project the reader has heard of (e.g., a national highway system; a major dam; a metropolitan water-and-sewer system; the original construction of a regional electrical grid; a national broadband or telecommunications buildout; a major airport hub; a continental rail network; a programme of public-housing construction). Trace, to the depth the public record supports: the inception decisions (who proposed; who decided; what political-economic forces shaped); the funding (who paid; how it was financed; what alternatives were displaced); the design (the engineering teams; the firms; the standards); the construction (labour; supply chain; impacted populations); the operation (operators; regulators; users); the long-term consequences (intended; unintended; ongoing); the contemporary debates (what is now contested; what the next generation of engineers may face). Deliverable: a 6-8 page document with the trace and a one-page reflection on what the trace reveals about engineering at civilization scale.

\textbf{Hazard.} The trace is necessarily incomplete; the public record on civilization-scale projects is large and partly contested. The point is the discipline of recognising the project's actual scope and the substantive engagement with what the engineering scope did not address; the document's value is in the substantive engagement, not in the trace's completeness.
\end{project}

\section{Exercises}

\subsection*{Calculation}\pathtag{core}

\begin{exercise}
A metropolitan water system serves 5 million people at an average use of $200\,\text{L/person/day}$. Compute the total daily delivery in m$^3$ and the annual delivery. Compare to the volume of a substantial reservoir (e.g., $10^9\,\text{m}^3$). Identify the implications for storage and source planning.
\end{exercise}

\begin{exercise}
A national infrastructure investment programme allocates 1\% of GDP per year for 20 years to renewal of aging infrastructure. For a country with annual GDP of \$2 trillion, compute the total investment. Compare to the ASCE-style backlog estimates (typically several trillion dollars for major economies).
\end{exercise}

\begin{exercise}
A civilization-scale project's design life is 100 years; its construction takes 10 years; the political cycle that funds it is 4 years. Compute the number of political cycles spanning the project's life; identify the cohort-thinking implication.
\end{exercise}

\subsection*{Derivation}\pathtag{standard}

\begin{exercise}
For a civilization-scale infrastructure decision with present cost $C$ and benefits accruing over $T$ years at $B$ per year, derive the present value at discount rate $r$. Identify the discount rate at which the present-value calculation favours the investment versus the rate at which it does not. Discuss the implications for long-horizon infrastructure decisions when market interest rates do not represent future generations' interests.
\end{exercise}

\begin{exercise}
For a system-of-systems with $n$ interconnected infrastructures, derive the conditions under which a failure in one infrastructure cascades to others (drawing on Volume X Chapter 7). Identify the design moves that reduce the cascade probability.
\end{exercise}

\subsection*{Estimation}\pathtag{standard}

\begin{exercise}
Estimate the engineering workforce required to maintain the existing infrastructure stock of a country with population $10^8$ at current condition. Identify the dominant categories (civil, electrical, mechanical, water, transportation) and compare to the actual workforce.
\end{exercise}

\begin{exercise}
Estimate the cost per displaced person of a major civilization-scale infrastructure project (a dam; a highway; an airport). Compare to the cost per beneficiary; identify the implications for the distributive analysis.
\end{exercise}

\subsection*{Simulation}\pathtag{standard}

\begin{exercise}
Implement a simulation of a multi-decade infrastructure investment programme under varying political-cycle volatility. Sweep the political-cycle parameter and identify the conditions under which the infrastructure can be maintained versus the conditions under which the cyclical funding produces accumulated deterioration.
\end{exercise}

\begin{exercise}
Simulate a system-of-systems with infrastructure dependencies. Model the failure cascade across infrastructures under various dependency strengths; identify the configurations that produce systemic-risk amplification versus the configurations that contain failures.
\end{exercise}

\subsection*{Design}\pathtag{standard}

\begin{exercise}
Design an institutional structure for long-horizon civilization-scale engineering decisions that explicitly represents future generations' interests. Specify the representation mechanism, the decision-rule's weighting of present versus future, the accountability for decisions whose consequences arrive after the decision-makers' departure.
\end{exercise}

\begin{exercise}
Design a process by which engineering projects' scope assessment includes identification of whether the project is civilization-scale in its consequences. Specify the criteria, the assessment procedure, the implications for the engineering scope, the institutional response when the assessment identifies civilization-scale implications.
\end{exercise}

\subsection*{Judgement}\pathtag{mastery}

\begin{exercise}
Argue: the engineering profession's responsibility extends to engaging with the political-economic context of its work; engineers who refuse the engagement on the grounds that "politics is not engineering" have misunderstood the profession at civilization scale. Argue the opposite. Identify the empirical evidence from the Volume X cases and from the civilization-scale projects of this chapter that bears on the dispute.
\end{exercise}

\begin{exercise}
Discuss the Aral Sea collapse as a civilization-scale engineering failure. Identify the engineering decisions that contributed; identify the institutional and political conditions; identify what would have been required of the engineering profession's discipline to surface the civilizational consequences before the project's execution. Discuss whether the profession would have had the institutional authority to act on the discipline.
\end{exercise}

\begin{exercise}
Argue: civilization-scale projects benefit from explicit decadal-to-multi-century planning; the absence of such planning produces the cumulative failures this chapter has documented. Identify the institutional changes that would support such planning and the political-economic constraints that resist it.
\end{exercise}

\begin{exercise}
Discuss the engineer's role in projects whose civilization-scale consequences the engineer believes are negative on net but the engineer's employer is committed to. Identify the engineer's options under the codes of ethics (Volume XI Chapter 11); discuss the institutional structures that would protect or fail to protect the engineer's professional judgment.
\end{exercise}
